% GENERATED FILE. Edit canonical lessons, then run npm run generate:books.
% canonical-source-hash: fnv1a64:19be94ac97b1c852
% canonical-lessons: JA-C366-saiaku, JA-C366-kesshin-suru, JA-C366-junjo, JA-C366-yuusen-suru, JA-C366-aseru, JA-R366-review-362-366

\chapter{If the Worst Happens}
\label{ch:ja-if-the-worst-happens}
\hlchaptermodality{366}

\begin{chapteropening}
\textbf{By the end of this chapter:} \emph{I can say what I would do if the worst happened.}

Complete the last lesson of chapter 366: I can say what I would do if the worst happened.
\end{chapteropening}

\section[\texorpdfstring{\ja{さいあく} (saiaku)}{さいあく (saiaku)}]{\ja{さいあく} (saiaku) — the worst}

\label{lesson:JA-C366-saiaku}



\begin{quote}
Before the new one: say the Japanese for a pension, then the Japanese for an average.
\end{quote}

\subsection*{You'll want to know: \ja{さいあく}}
\textbf{\ja{さいあく}} — \emph{saiaku} — \textquotedblleft{}the worst\textquotedblright{}.

\begin{quote}
\textbf{\ja{さいあく、でんしゃが} \ja{とまったら、あるいて} \ja{かえります。}} — \emph{saiaku, densha ga tomattara, aruite kaerimasu} — \textquotedblleft{}If the worst comes to the worst and the trains stop, I will walk home.\textquotedblright{}
\end{quote}

A na-adjective: \ja{な} before a noun.

\begin{cousinweb}
\textbf{\ja{さいあく}} (\emph{saiaku}) is Sino-Japanese (kango): two Chinese characters read with their on'yomi, readings borrowed from Chinese. They are \emph{sai}, \textquotedblleft{}utmost\textquotedblright{} (Middle Chinese \emph{tswajH}) and \emph{aku}, \textquotedblleft{}bad\textquotedblright{} (\emph{'ak}) (KANJIDIC2; Middle Chinese in Baxter's transcription).
\end{cousinweb}

\begin{grammarlens}[title={What you've built}]
One more word for saying what happens if.
\end{grammarlens}

\subsection*{Your turn}
\begin{itemize}
\raggedright
  \item \emph{Say it:} \emph{saiaku}
  \item \emph{Say it:} \emph{saiaku}, once more
  \item \emph{Say it:} \emph{saiaku, densha ga tomattara, aruite kaerimasu} — If the worst comes to the worst and the trains stop, I will walk home.
  \item \emph{Recall:} say the Japanese for a pension, then the Japanese for an average, then say \emph{saiaku} again
\end{itemize}

\begin{tcolorbox}[breakable,colback=teal!4,colframe=teal!35!black,title={Before you move on}]
What does \ja{さいあく} mean? (\textquotedblleft{}The worst\textquotedblright{}.) Say it once more.
\end{tcolorbox}
\section[\texorpdfstring{\ja{けっしんする} (kesshin suru)}{けっしんする (kesshin suru)}]{\ja{けっしんする} (kesshin suru) — to make up one's mind}

\label{lesson:JA-C366-kesshin-suru}



\begin{quote}
Before the new one: say the Japanese for an average, then the Japanese for the worst.
\end{quote}

\subsection*{You'll want to know: \ja{けっしんする}}
\textbf{\ja{けっしんする}} — \emph{kesshin suru} — \textquotedblleft{}to make up one's mind\textquotedblright{}.

\begin{quote}
\textbf{\ja{ははが} \ja{さんせいしたら、けっしんします。}} — \emph{haha ga sansei shitara, kesshin shimasu} — \textquotedblleft{}If my mother agrees, I will make up my mind.\textquotedblright{}
\end{quote}

Polite \textbf{\ja{けっしんします}}, \emph{kesshinshimasu}; a \ja{する} verb. Past (polite): \textbf{\ja{あには} \ja{けっしんしました。}}, \emph{ani wa kesshin shimashita}.

\begin{cousinweb}
\textbf{\ja{けっしんする}} (\emph{kesshin suru}) is built from a Sino-Japanese noun and the native verb \emph{suru}, \textquotedblleft{}to do\textquotedblright{}. The noun, \emph{kesshin}, is two Chinese characters read with their on'yomi, readings borrowed from Chinese: \emph{ketsu}, \textquotedblleft{}decide\textquotedblright{} (Middle Chinese \emph{kwet}) and \emph{shin}, \textquotedblleft{}heart\textquotedblright{} (\emph{sim}) (KANJIDIC2; JMdict; Middle Chinese in Baxter's transcription). The two join as \emph{kes-shin}.
\end{cousinweb}

\begin{grammarlens}[title={What you've built}]
One more word for saying what happens if.
\end{grammarlens}

\subsection*{Your turn}
\begin{itemize}
\raggedright
  \item \emph{Say it:} \emph{kesshin suru}
  \item \emph{Say it:} \emph{kesshin suru}, once more
  \item \emph{Say it:} \emph{haha ga sansei shitara, kesshin shimasu} — If my mother agrees, I will make up my mind.
  \item \emph{Recall:} say the Japanese for an average, then the Japanese for the worst, then say \emph{kesshin suru} again
\end{itemize}

\begin{tcolorbox}[breakable,colback=teal!4,colframe=teal!35!black,title={Before you move on}]
What does \ja{けっしんする} mean? (\textquotedblleft{}To make up one's mind\textquotedblright{}.) Say it once more.
\end{tcolorbox}
\section[\texorpdfstring{\ja{じゅんじょ} (junjo)}{じゅんじょ (junjo)}]{\ja{じゅんじょ} (junjo) — order; sequence}

\label{lesson:JA-C366-junjo}



\begin{quote}
Before the new one: say the Japanese for the worst, then the Japanese for to make up one's mind.
\end{quote}

\subsection*{You'll want to know: \ja{じゅんじょ}}
\textbf{\ja{じゅんじょ}} — \emph{junjo} — \textquotedblleft{}order; sequence\textquotedblright{}.

\begin{quote}
\textbf{\ja{じゅんじょを} \ja{まちがえたら、うごきません。}} — \emph{junjo o machigaetara, ugokimasen} — \textquotedblleft{}If you do it in the wrong order, it won't work.\textquotedblright{}
\end{quote}

\begin{cousinweb}
\textbf{\ja{じゅんじょ}} (\emph{junjo}) is Sino-Japanese (kango): two Chinese characters read with their on'yomi, readings borrowed from Chinese. They are \emph{jun}, \textquotedblleft{}order\textquotedblright{} (Middle Chinese \emph{zywinH}) and \emph{jo}, \textquotedblleft{}order\textquotedblright{} (\emph{zjoX}) (KANJIDIC2; Middle Chinese in Baxter's transcription).
\end{cousinweb}

\begin{grammarlens}[title={What you've built}]
One more word for saying what happens if.
\end{grammarlens}

\subsection*{Your turn}
\begin{itemize}
\raggedright
  \item \emph{Say it:} \emph{junjo}
  \item \emph{Say it:} \emph{junjo}, once more
  \item \emph{Say it:} \emph{junjo o machigaetara, ugokimasen} — If you do it in the wrong order, it won't work.
  \item \emph{Recall:} say the Japanese for the worst, then the Japanese for to make up one's mind, then say \emph{junjo} again
\end{itemize}

\begin{tcolorbox}[breakable,colback=teal!4,colframe=teal!35!black,title={Before you move on}]
What does \ja{じゅんじょ} mean? (\textquotedblleft{}Order; sequence\textquotedblright{}.) Say it once more.
\end{tcolorbox}
\section[\texorpdfstring{\ja{ゆうせんする} (yūsen suru)}{ゆうせんする (yūsen suru)}]{\ja{ゆうせんする} (yūsen suru) — to put first; to give priority to}

\label{lesson:JA-C366-yuusen-suru}



\begin{quote}
Before the new one: say the Japanese for to make up one's mind, then the Japanese for the sequence.
\end{quote}

\subsection*{You'll want to know: \ja{ゆうせんする}}
\textbf{\ja{ゆうせんする}} — \emph{yūsen suru} — \textquotedblleft{}to put first; to give priority to\textquotedblright{}.

\begin{quote}
\textbf{\ja{しけんが} \ja{あったら、べんきょうを} \ja{ゆうせんします。}} — \emph{shiken ga attara, benkyō o yūsen shimasu} — \textquotedblleft{}If there is an exam, I put studying first.\textquotedblright{}
\end{quote}

Polite \textbf{\ja{ゆうせんします}}, \emph{yūsenshimasu}; a \ja{する} verb. Past (polite): \textbf{\ja{ははは} \ja{かぞくを} \ja{ゆうせんしました。}}, \emph{haha wa kazoku o yūsen shimashita}.

\begin{cousinweb}
\textbf{\ja{ゆうせんする}} (\emph{yūsen suru}) is built from a Sino-Japanese noun and the native verb \emph{suru}, \textquotedblleft{}to do\textquotedblright{}. The noun, \emph{yūsen}, is two Chinese characters read with their on'yomi, readings borrowed from Chinese: \emph{yū}, \textquotedblleft{}excel\textquotedblright{} (Middle Chinese \emph{'juw}) and \emph{sen}, \textquotedblleft{}ahead\textquotedblright{} (\emph{sen}) (KANJIDIC2; JMdict; Middle Chinese in Baxter's transcription).
\end{cousinweb}

\begin{grammarlens}[title={What you've built}]
One more word for saying what happens if.
\end{grammarlens}

\subsection*{Your turn}
\begin{itemize}
\raggedright
  \item \emph{Say it:} \emph{yūsen suru}
  \item \emph{Say it:} \emph{yūsen suru}, once more
  \item \emph{Say it:} \emph{shiken ga attara, benkyō o yūsen shimasu} — If there is an exam, I put studying first.
  \item \emph{Recall:} say the Japanese for to make up one's mind, then the Japanese for the sequence, then say \emph{yūsen suru} again
\end{itemize}

\begin{tcolorbox}[breakable,colback=teal!4,colframe=teal!35!black,title={Before you move on}]
What does \ja{ゆうせんする} mean? (\textquotedblleft{}To put first; to give priority to\textquotedblright{}.) Say it once more.
\end{tcolorbox}
\section[\texorpdfstring{\ja{あせる} (aseru)}{あせる (aseru)}]{\ja{あせる} (aseru) — to be impatient; to be in a rush}

\label{lesson:JA-C366-aseru}



\begin{quote}
Before the new one: say the Japanese for the sequence, then the Japanese for to put first.
\end{quote}

\subsection*{You'll want to know: \ja{あせる}}
\textbf{\ja{あせる}} — \emph{aseru} — \textquotedblleft{}to be impatient; to be in a rush\textquotedblright{}.

\begin{quote}
\textbf{\ja{あせったら、まちがえるでしょう。}} — \emph{asettara, machigaeru deshō} — \textquotedblleft{}If you rush, you will probably make mistakes.\textquotedblright{}
\end{quote}

Polite \textbf{\ja{あせります}}, \emph{aserimasu}; a godan verb. Past (polite): \textbf{\ja{おとうとが} \ja{あせりました。}}, \emph{otōto ga aserimashita}.

\begin{cousinweb}
\textbf{\ja{あせる}} (\emph{aseru}) is a native Japanese word (wago), inherited rather than borrowed from Chinese, and its older origin is not settled. Its kanji was attached to it later as a kun'yomi, a native reading given to a Chinese character (JMdict, KANJIDIC2).
\end{cousinweb}

\begin{grammarlens}[title={What you've built}]
One more word for saying what happens if.
\end{grammarlens}

\subsection*{Your turn}
\begin{itemize}
\raggedright
  \item \emph{Say it:} \emph{aseru}
  \item \emph{Say it:} \emph{aseru}, once more
  \item \emph{Say it:} \emph{asettara, machigaeru deshō} — If you rush, you will probably make mistakes.
  \item \emph{Recall:} say the Japanese for the sequence, then the Japanese for to put first, then say \emph{aseru} again
\end{itemize}

\begin{tcolorbox}[breakable,colback=teal!4,colframe=teal!35!black,title={Before you move on}]
What does \ja{あせる} mean? (\textquotedblleft{}To be impatient; to be in a rush\textquotedblright{}.) Say it once more.
\end{tcolorbox}
\section[(dialogue)]{Looking back: from What Will It Cost? to If the Worst Happens}

\label{lesson:JA-R366-review-362-366}



\begin{quote}
Say the Japanese for an expense and to calculate.
\end{quote}

\subsection*{Your turn}
Say these aloud:

\begin{itemize}
\raggedright
  \item an expense, to calculate, the total, prices in general, a rental deposit
  \item the bill, the payment, consumption tax, half price, luxury
  \item damage to property, to pay for damage, to pay taxes, to be in charge of, responsibility
  \item the economy, business conditions, an investment, a pension, an average
  \item the worst, to make up one's mind, the sequence, to put first, to be impatient
\end{itemize}

\begin{tcolorbox}[breakable,colback=teal!4,colframe=teal!35!black,title={Before you move on}]
An expense? (\textbf{\ja{ひよう}}, \emph{hiyō}.) To calculate? (\textbf{\ja{けいさんする}}, \emph{keisan suru}.) The total? (\textbf{\ja{ごうけい}}, \emph{gōkei}.) Prices in general? (\textbf{\ja{ぶっか}}, \emph{bukka}.) A rental deposit? (\textbf{\ja{しききん}}, \emph{shikikin}.) The bill? (\textbf{\ja{かいけい}}, \emph{kaikei}.) The payment? (\textbf{\ja{だいきん}}, \emph{daikin}.) Consumption tax? (\textbf{\ja{しょうひぜい}}, \emph{shōhizei}.) Half price? (\textbf{\ja{はんがく}}, \emph{hangaku}.) Luxury? (\textbf{\ja{ぜいたく}}, \emph{zeitaku}.) Damage to property? (\textbf{\ja{そんがい}}, \emph{songai}.) To pay for damage? (\textbf{\ja{べんしょうする}}, \emph{benshō suru}.) To pay taxes? (\textbf{\ja{おさめる}}, \emph{osameru}.) To be in charge of? (\textbf{\ja{うけもつ}}, \emph{ukemotsu}.) Responsibility? (\textbf{\ja{せきにん}}, \emph{sekinin}.) The economy? (\textbf{\ja{けいざい}}, \emph{keizai}.) Business conditions? (\textbf{\ja{けいき}}, \emph{keiki}.) An investment? (\textbf{\ja{とうし}}, \emph{tōshi}.) A pension? (\textbf{\ja{ねんきん}}, \emph{nenkin}.) An average? (\textbf{\ja{へいきん}}, \emph{heikin}.) The worst? (\textbf{\ja{さいあく}}, \emph{saiaku}.) To make up one's mind? (\textbf{\ja{けっしんする}}, \emph{kesshin suru}.) The sequence? (\textbf{\ja{じゅんじょ}}, \emph{junjo}.) To put first? (\textbf{\ja{ゆうせんする}}, \emph{yūsen suru}.) To be impatient? (\textbf{\ja{あせる}}, \emph{aseru}.)
\end{tcolorbox}
